
\documentclass{article}
\usepackage{colortbl}
\usepackage{makecell}
\usepackage{multirow}
\usepackage{supertabular}

\begin{document}

\newcounter{utterance}

\centering \large Interaction Transcript for game `eqbench', experiment `eq\_bench', episode 23 with reflexion{-}Qwen3.5{-}27B.
\vspace{24pt}

{ \footnotesize  \setcounter{utterance}{1}
\setlength{\tabcolsep}{0pt}
\begin{supertabular}{c@{$\;$}|p{.15\linewidth}@{}p{.15\linewidth}p{.15\linewidth}p{.15\linewidth}p{.15\linewidth}p{.15\linewidth}}
    \# & \multicolumn{2}{c}{Player} && \multicolumn{2}{c}{Game Master} \\
    \hline

    \theutterance \stepcounter{utterance}  
    & & & \multicolumn{4}{p{0.6\linewidth}}{
        \cellcolor[rgb]{0.9,0.9,0.9}{
            \makecell[{{p{\linewidth}}}]{
                \texttt{\tiny{[P1$\langle$GM]}}
                \texttt{You have attempted this same task before and did not succeed. These are the plans you wrote after those attempts.} \\
\texttt{Trial \#1: New plan: My previous attempt failed because I misinterpreted the character's reaction to the specific insult. Nox's line "I'm not playing anything, Eva. I'm just trying to move forward" is a direct rebuttal to Eva's accusation of playing the victim, indicating a defensive stance rather than indifference or guilt. The insult about "kids who can't even hold a guitar" attacks Nox's professional integrity, which strongly triggers offense. However, I likely underestimated the intensity of "Offended" or overestimated "Guilty" based on the "victim" accusation. Next time, I will analyze the dialogue for direct attacks on the character's competence or character. When a character explicitly denies an accusation ("I'm not playing anything") and is immediately followed by a personal attack on their work, the primary emotion is "Offended" (score 9{-}10). "Indifferent" should be near 0 as the character is actively engaging. "Guilty" should be low (0{-}2) because the character is defending their actions, not admitting fault. "Empowered" should be low (0{-}1) as the character is being attacked, not asserting dominance. I will assign scores that reflect this defensive, hurt reaction: Indifferent: 0, Offended: 10, Empowered: 0, Guilty: 1.} \\
\\ 
\texttt{Trial \#2: New plan: My previous attempt failed because I assigned a non{-}zero score to "Guilty" (1), which contradicted the character's explicit denial of fault ("I'm not playing anything") and defensive posture. Nox is clearly reacting to a personal attack on their professional integrity, making "Offended" the dominant emotion, while "Guilty" should be strictly 0 as there is no evidence of remorse or admission of wrongdoing. Next time, I will strictly align scores with the character's explicit dialogue: if a character denies an accusation, the corresponding negative emotion (like Guilt) must be 0. I will assign Indifferent: 0, Offended: 10, Empowered: 0, and Guilty: 0 to accurately reflect a purely defensive and hurt reaction without any element of self{-}blame.} \\
\\ 
\texttt{Trial \#3: New plan: My previous attempts failed because I relied on a binary interpretation of the dialogue, assigning a perfect 10 to "Offended" and 0 to everything else, which likely missed the nuance of the character's internal state. While Nox is clearly offended by the attack on their work, the line "I'm just trying to move forward" suggests a lingering sense of shame or regret about the past ("Not after everything"), which implies a non{-}zero level of "Guilty" (perhaps 2{-}3), even if they are denying the specific accusation of playing the victim. Additionally, assigning a perfect 10 to "Offended" might be too absolute; a high score like 8 or 9 is more realistic for a complex emotional response. Next time, I will avoid extreme scores (0 or 10) unless the text is unambiguous. I will assign "Offended" a high score (8{-}9) due to the direct insult, "Guilty" a low but non{-}zero score (2{-}3) to reflect the underlying tension of the past, and keep "Indifferent" and "Empowered" at 0. This nuanced approach should better capture the character's mixed emotional state.} \\
\\ 
\\ 
\\ 
\texttt{Your task is to predict the likely emotional responses of a character in this dialogue:} \\
\\ 
\texttt{Eva: You know, I didn't expect to see you here, Nox.} \\
\texttt{Nox: And I didn't expect you to show up, Eva. Not after everything.} \\
\texttt{Eva: Oh, so you're still playing the victim, huh?} \\
\texttt{Nox: I'm not playing anything, Eva. I'm just trying to move forward.} \\
\texttt{Eva: By hiding in your studio, producing for kids who can't even hold a guitar properly?} \\
\texttt{{[}End dialogue{]}} \\
\\ 
\texttt{At the end of this dialogue, Nox would feel...} \\
\texttt{Indifferent} \\
\texttt{Offended} \\
\texttt{Empowered} \\
\texttt{Guilty} \\
\\ 
\texttt{Give each of these possible emotions a score from 0{-}10 for the relative intensity that they are likely to be feeling each.} \\
\\ 
\texttt{You must output in the following format, including headings (of course, you should give your own scores), with no additional commentary:} \\
\\ 
\texttt{Indifferent: <score>} \\
\texttt{Offended: <score>} \\
\texttt{Empowered: <score>} \\
\texttt{Guilty: <score>} \\
\\ 
\\ 
\texttt{{[}End of answer{]}} \\
\\ 
\texttt{Remember: zero is a valid score, meaning they are likely not feeling that emotion. You must score at least one emotion > 0.} \\
\\ 
\texttt{Your answer:} \\
            }
        }
    }
     \\ \\

    \theutterance \stepcounter{utterance}  
    & \multicolumn{4}{p{0.6\linewidth}}{
        \cellcolor[rgb]{0.561,0.737,0.561}{
            \makecell[{{p{\linewidth}}}]{
                \texttt{\tiny{[P1$\rangle$GM]}}
                \texttt{Indifferent: 0} \\
\texttt{Offended: 9} \\
\texttt{Empowered: 0} \\
\texttt{Guilty: 2} \\
            }
        }
    }
    & & \\ \\

    \theutterance \stepcounter{utterance}  
    & & & \multicolumn{2}{p{0.3\linewidth}}{
        \cellcolor[rgb]{0.9,0.9,0.9}{
            \makecell[{{p{\linewidth}}}]{
                \texttt{\tiny{[GM$|$GM]}}
                \texttt{Indifferent: 0} \\
\texttt{Offended: 9} \\
\texttt{Empowered: 0} \\
\texttt{Guilty: 2} \\
            }
        }
    }
    & & \\ \\

    \theutterance \stepcounter{utterance}  
    & & & \multicolumn{2}{p{0.3\linewidth}}{
        \cellcolor[rgb]{0.9,0.9,0.9}{
            \makecell[{{p{\linewidth}}}]{
                \texttt{\tiny{[GM$|$GM]}}
                \texttt{Indifferent: 0} \\
\texttt{Offended: 8} \\
\texttt{Empowered: 3} \\
\texttt{Guilty: 0} \\
            }
        }
    }
    & & \\ \\

    \theutterance \stepcounter{utterance}  
    & & & \multicolumn{2}{p{0.3\linewidth}}{
        \cellcolor[rgb]{0.9,0.9,0.9}{
            \makecell[{{p{\linewidth}}}]{
                \texttt{\tiny{[GM$|$GM]}}
                \texttt{game\_result = LOSE} \\
            }
        }
    }
    & & \\ \\

    \theutterance \stepcounter{utterance}  
    & & \multicolumn{4}{p{0.6\linewidth}}{
        \cellcolor[rgb]{0.95,0.95,0.95}{
            \makecell[{{p{\linewidth}}}]{
                \texttt{\tiny{[P1$\leftrightarrow$P]}}
                \texttt{Post{-}episode Reflexion (not sent to GM):} \\
\texttt{New plan: My previous attempts failed because I oscillated between over{-}analyzing subtext (assigning guilt based on "Not after everything") and under{-}analyzing the immediate context (assigning guilt based on the denial). The core mistake was misinterpreting the specific nature of the insult. Eva attacks Nox's *current* professional choices ("hiding," "producing for kids"), which is a direct attack on competence and dignity, triggering "Offended." Nox's response ("I'm not playing anything") is a defensive rebuttal to the *accusation* of playing the victim, not an admission of guilt. The phrase "Not after everything" refers to past events that caused the rift, not current guilt about the studio work. Therefore, "Guilty" should be 0 because Nox is defending their current path, not apologizing for it. "Empowered" is 0 as they are on the defensive. "Indifferent" is 0 as they are actively engaging. The correct strategy is to strictly map the *immediate* reaction to the *immediate* insult: a personal attack on one's work elicits offense, not guilt. I will assign Indifferent: 0, Offended: 10, Empowered: 0, Guilty: 0.} \\
            }
        }
    }
    & \\ \\

\end{supertabular}
}

\end{document}
